\section*{Hoofdstuk \ref{ch:b3:genomics} --- Genomica en sequencing}

\begin{solution}{exo:b3:genomics:1}
Een \omterm{def:b3:genomics:sanger}{dideoxynucleotide} heeft geen 3$'$-hydroxylgroep, dus kan er geen
fosfodiesterbinding met het volgende nucleotide worden gemaakt en stopt
de keten. Met alleen dideoxyvormen zou elke keten bij de eerste positie
stoppen; met alleen normale vormen zou er geen enkele stoppen. Het
mengsel maakt de terminatie op elke positie een toevalsgebeurtenis,
zodat de producten een volledige ladder vormen, één lengte voor elke
base van de matrijs.
\end{solution}

\begin{solution}{exo:b3:genomics:2}
$4\times 10^{8}\times \qty{300}{bp} = 1.2\times 10^{11}$ basen. Mens:
$1.2\times 10^{11}/3.2\times 10^{9} \approx 38\times$. Bacterie:
$1.2\times 10^{11}/5\times 10^{6} = \num{24000}\times$.
\end{solution}

\begin{solution}{exo:b3:genomics:3}
Een contig is een aaneengesloten sequentie die uit overlappende \omterm{def:b3:genomics:ngs}{reads} is
geassembleerd; een \omterm{def:b3:genomics:assembly}{scaffold} is een geordende, gerichte reeks contigs met
gaten van geschatte grootte ertussen; de \omterm{def:b3:genomics:assembly}{N50} is de contiglengte waarbij
de gesorteerde contigs de helft van de geassembleerde basen bereiken.
Hier bevatten de tien contigs van \qty{8}{Mb} al \qty{80}{Mb}, voorbij
het halverwegepunt van \qty{50}{Mb} (de zevende bereikt \qty{56}{Mb}):
\omterm{def:b3:genomics:assembly}{N50} $= \qty{8}{Mb}$.
\end{solution}

\begin{solution}{exo:b3:genomics:4}
De \omterm{prop:b3:genomics:sizes}{genoomgrootte} volgt de complexiteit van een organisme niet: een ui
heeft vijf keer zoveel DNA als een mens, een longvis veertig keer;
gist en de worm verschillen tienvoudig in DNA bij vergelijkbare
aantallen genen. Het overschot is niet-coderend: verplaatsbare
elementen en hun resten, introns, satellietherhalingen.
\end{solution}

\begin{solution}{exo:b3:genomics:5}
$e^{-c} = 10^{-6}$ geeft $c = 6\ln 10 = 13.8$. Voor $c = 8$: $N =
cG/L = 8\times 5\times 10^{6}/150 \approx \num{267000}$ \omterm{def:b3:genomics:ngs}{reads}, $\theta
= 30/150 = 0.2$, contigs $= N e^{-6.4} = \num{267000}\times 0.00166
\approx 440$.
\end{solution}

\begin{solution}{exo:b3:genomics:6}
$P(X < 8) = P(X \le 7) \approx P\bigl(Z < (7.5 - 15)/2.74\bigr) =
P(Z < -2.74) \approx 0.003$. Bij $30\times$ worden beide allelen van
een heterozygoot vrijwel altijd vele keren gezien, is de dekking
ongelijkmatig (GC-rijke gebieden krijgen minder \omterm{def:b3:genomics:ngs}{reads}, zodat sommige
plaatsen de helft van het gemiddelde zien), en blijven er genoeg \omterm{def:b3:genomics:ngs}{reads}
over om een echt allel van de fouten van $10^{-3}$ te scheiden.
\end{solution}

\begin{solution}{exo:b3:genomics:7}
Elke \omterm{def:b3:genomics:ngs}{read} van \qty{150}{bp} uit de herhaling is identiek, van welke
kopie hij ook komt, dus bouwt de assembler één knoop van \qty{6}{kb}
met $500$ paden die binnenkomen en $500$ die vertrekken; hij kan niet
zien welke ingang bij welke uitgang hoort, en de \omterm{def:b3:genomics:assembly}{assemblage} breekt in
$500$ gaten, met de herhaling één keer aanwezig bij een dekking van
$500$ keer het gemiddelde. \omterm{def:b3:genomics:ngs}{Reads} die langer zijn dan de herhaling plus
unieke flanken aan beide zijden --- \qty{8}{kb} of meer --- overspannen
elke kopie en lossen haar op.
\end{solution}

\begin{solution}{exo:b3:genomics:8}
$4.8\times 10^{7}/1000 = \num{48000}$ coderende verschillen; een derde
daarvan, ongeveer \num{16000}, verandert het eiwit.
\end{solution}

\begin{solution}{exo:b3:genomics:9}
Verwachte valse positieven $10^{6}\times 5\times 10^{-8} = 0.05$. Een
relatief risico van $1.15$ op een ziekte van \qty{2}{\%} verandert
enkele gevallen per duizend; duizend mensen bevatten ongeveer twintig
gevallen, veel te weinig om het te zien (het onderscheidend vermogen
groeit met het aantal gevallen en met het kwadraat van het effect). Voor
een individu is $0.3$ procentpunt betekenisloos. Maar het allel
markeert een gen waarvan een bescheiden verandering in activiteit het
ziekterisico verandert --- het gen, en zijn route, kunnen een
aangrijpingspunt voor een geneesmiddel zijn waarvan volledige remming
een groot effect heeft.
\end{solution}

\begin{solution}{exo:b3:genomics:10}
Populatiegenetisch: bij bacteriën, met enorme populaties, verwijdert de
selectie zelfs licht kostbare inserties; bij zoogdieren, met kleine
effectieve populaties, laat de drift licht schadelijk niet-coderend DNA
zich ophopen --- gesteund door de omgekeerde samenhang tussen
\omterm{prop:b3:genomics:sizes}{genoomgrootte} en populatiegrootte over afstammingslijnen heen,
ondergraven door uitzonderingen. Structureel: eukaryote genomen worden
binnengedrongen door transposons die bacteriën, met een voorkeur voor
deletie en zonder meiotisch toevluchtsoord voor zelfzuchtige elementen,
opruimen --- gesteund door de samenhang tussen \omterm{prop:b3:genomics:sizes}{genoomgrootte} en
transposoninhoud. Regulerend: complexe ontwikkeling vergt meer
regulerend DNA --- gesteund door de geconserveerde niet-coderende
elementen rond ontwikkelingsgenen, ondergraven door het feit dat maar
\qty{8}{\%} van het menselijk genoom enige selectie vertoont, zodat het
meeste niet-coderende DNA niet regulerend is.
\end{solution}

\begin{solution}{exo:b3:genomics:11}
Beginpunten van beide soorten vallen met totale dichtheid $\rho = (N_{1} +
N_{2})/G$. Een \omterm{def:b3:genomics:ngs}{read} van lengte $L_{i}$ is de meest rechtse van zijn
contig als er geen \omterm{def:b3:genomics:ngs}{read} van welke soort dan ook begint in de
$L_{i} - T$ basen erna: kans $e^{-\rho(L_{i} - T)}$. Dus contigs $= N_{1}
e^{-\rho(L_{1} - T)} + N_{2} e^{-\rho(L_{2} - T)}$. Een lange \omterm{def:b3:genomics:ngs}{read} is de
meest rechtse met een exponentieel kleinere kans dan een korte, en elke
lange \omterm{def:b3:genomics:ngs}{read} haalt ook over zijn hele lengte de kans op een gat weg;
hetzelfde aantal basen in korte \omterm{def:b3:genomics:ngs}{reads} voegt vele beginpunten toe maar
elk beschermd venster is kort, zodat de lange \omterm{def:b3:genomics:ngs}{reads} winnen.
\end{solution}

\begin{solution}{exo:b3:genomics:12}
Twee \omterm{def:b3:genomics:comparative}{volledige-genoomduplicaties} voorspellen dat het viervoudige patroon
genoombreed is: viertallen paraloge chromosoomsegmenten (paralogons)
die dezelfde genfamilies in dezelfde volgorde dragen, met duplicaties
die voor alle families op dezelfde tijd worden gedateerd; één enkel
cluster bij amphioxus en \emph{Ciona}, die vóór de gebeurtenissen zijn
afgesplitst; en bij de prikken, die er rond zijn afgesplitst, een
tussenvorm of een onafhankelijk afgeleide toestand. Onafhankelijke
duplicaties van alleen het Hox-cluster voorspellen paralogons voor
alleen Hox, buren met een andere geschiedenis, en dateringen van de
duplicatie die per familie verschillen. De genoombrede paralogons en de
gedeelde datering, zoals waargenomen, pleiten voor
\omterm{def:b3:genomics:comparative}{volledige-genoomduplicatie}.
\end{solution}

\begin{solution}{pb:b3:genomics:1}
\textbf{1.} $0.002\times 6\times 10^{8} = 1.2\times 10^{6}$ \omterm{def:b3:genomics:ngs}{reads};
$c = 1.2\times 10^{6}\times 150/4.6\times 10^{6} \approx 39$.
\textbf{2.} $e^{-39} \approx 10^{-17}$: feitelijk geen enkele base
onbedekt.
\textbf{3.} $\theta = 0.2$; contigs $= 1.2\times 10^{6} e^{-31} \approx
0$: in theorie één contig, met gaten alleen bij herhalingen.
\textbf{4.} $c = \num{30000}\times 150/4.6\times 10^{6} = 0.98$;
onbedekt $e^{-0.98} = 0.38$, ongeveer \qty{1.7}{Mb}; contigs $= \num{30000}
\times e^{-0.78} \approx \num{13700}$.
\textbf{5.} De zeven operons geven identieke \omterm{def:b3:genomics:ngs}{reads} en klappen samen tot
één contig van \qty{5}{kb}, binnengegaan door zeven unieke linkerflanken
en verlaten door zeven rechterflanken waarvan de koppeling onbekend is:
$14$ contiguiteinden, zeven gaten.
\textbf{6.} Ongeveer \num{4600} genen; coderend $0.88\times 4.6\times 10^{6}
= 4.0\times 10^{6}$ bp, gemiddeld gen \qty{880}{bp}.
\textbf{7.} $0.998\times 6\times 10^{8}\times 150/3.2\times 10^{9}
\approx 28\times$.
\textbf{8.} Ongeveer $14$ \omterm{def:b3:genomics:ngs}{reads} per allel.
\textbf{9.} $28\times 10^{-3}/3 \approx 0.009$ \omterm{def:b3:genomics:ngs}{reads} die een gegeven
verkeerde base tonen --- de kans op drie of meer is ongeveer $10^{-7}$
--- en \qty{20}{\%} van $28$ is $5.6$ \omterm{def:b3:genomics:ngs}{reads}; een fout doorstaat geen van
beide toetsen, terwijl een echt heterozygoot allel in $14$ \omterm{def:b3:genomics:ngs}{reads} wordt
verwacht.
\textbf{10.} $3.2\times 10^{9}/1500 \approx 2.1\times 10^{6}$
heterozygote plaatsen; in coderende sequentie $4.8\times 10^{7}/1500 \approx
\num{32000}$.
\textbf{11.} Een \omterm{def:b3:genomics:ngs}{read} van \qty{150}{bp} uit een Alu past vrijwel even
goed op veel van de $1.1$ miljoen kopieën, dus wordt hij op de verkeerde
kopie gelegd of met weinig vertrouwen gekarteerd. Verschillen tussen
kopieën lijken dan op heterozygote varianten, en echte varianten gaan
daartussen schuil: variantbepalingen binnen Alu worden meestal
weggefilterd, en de elementen zijn blinde vlekken van sequencing met
korte \omterm{def:b3:genomics:ngs}{reads}.
\textbf{12.} $1.2\times 10^{-8}\times 6.4\times 10^{9} \approx 77$
nieuwe mutaties. Ongeveer $14$ \omterm{def:b3:genomics:ngs}{reads} zouden de variant bij het kind
moeten tonen; maar als de plaats bij een ouder door maar vijf \omterm{def:b3:genomics:ngs}{reads}
wordt bedekt, wordt een heterozygoot allel met kans $2^{-5} = 3\,\%$
gemist en wordt een geërfde variant ten onrechte als nieuw bestempeld
--- dus moeten de ouders even diep worden gesequenced als het kind.
\textbf{13.} $4.8\times 10^{7}\times 100/150 = 3.2\times 10^{7}$ \omterm{def:b3:genomics:ngs}{reads},
twintig keer minder dan het genoom: de sequencingkosten dalen
twintigvoudig, meer dan de kosten van de vangststap.
\textbf{14.} $1.1\times 10^{6}\times 300 = 3.3\times 10^{8}$ bp, ongeveer
\qty{10}{\%} van het genoom; \qty{10}{\%} van de \omterm{def:b3:genomics:ngs}{reads}, zo'n $6\times
10^{7}$.
\textbf{15.} $\num{60000}\times 1500 = 9\times 10^{7}$ bp; $9\times
10^{7}/1.6\times 10^{10} = 0.56\,\%$.
\textbf{16.} $0.012\times 2.9\times 10^{9} = 3.5\times 10^{7}$
substituties, $1.7\times 10^{7}$ per afstammingslijn; snelheid $1.7\times
10^{7}/(2.9\times 10^{9}\times 6.5\times 10^{6}) = 9\times 10^{-10}$
per base per jaar, $2.3\times 10^{-8}$ per generatie van 25 jaar ---
ongeveer tweemaal de stamboomsnelheid van $1.2\times 10^{-8}$, een
bekende discrepantie die op langere generatietijden in het verleden of
op een oudere splitsing wijst.
\textbf{17.} Vier kopieën van \num{16000} zouden \num{64000} zijn;
\num{20000} blijven over, dus van de \num{48000} extra kopieën
overleven er maar \num{4000}: \qty{92}{\%} van de duplicaten ging
verloren.
\textbf{18.} De worm en de mens hebben evenveel genen; de complexiteit
ligt in hoe ze worden gebruikt --- alternatieve splicing vermenigvuldigt
één gen tot duizenden eiwitten (\num{38016} voor \emph{Dscam}), de
regulatie combineert transcriptiefactoren en enhancers op
celtypespecifieke wijze, niet-coderende RNA's voegen lagen toe, en
eiwitten werken in netwerken samen.
\textbf{19.} Regulerende elementen (enhancers, promotoren, isolatoren),
genen voor \omterm{def:b3:rna-regulation:ncrna}{niet-coderend RNA}, signalen voor splitsing en lokalisatie,
replicatie-oorsprongen en structurele sequenties. Toets een
kandidaat-enhancer door hem vóór een reportergen in een transgeen embryo
te plaatsen en te kijken waar het reportergen tot expressie komt, en
verwijder het element daarna met CRISPR uit het genoom en meet de
naburige genen.
\textbf{20.} $0.05/10^{6} = 5\times 10^{-8}$; bij $p < 10^{-5}$ zijn
$10^{6}
\times 10^{-5} = 10$ valse positieven te verwachten.
\textbf{21.} Odds voor een niet-drager $0.009/0.991 = 0.00908$; bij één
kopie odds $0.00954$, risico $0.945\,\%$; bij twee kopieën odds
$0.01001$, risico $0.99\,\%$.
\textbf{22.} De score telt over de $200$ allelen het aantal
risicokopieën dat iemand draagt, gewogen met de log-oddsratio van elk
allel. De telling heeft gemiddelde $200\times 2\times 0.3 = 120$ en
standaardafwijking $\sqrt{200\times 2\times 0.3\times 0.7} \approx 9$;
de bovenste \qty{1}{\%} draagt ongeveer $21$ allelen meer dan gemiddeld,
en $1.05^{21} \approx
2.8$: bijna drievoudige odds ten opzichte van een gemiddeld persoon, uit
allelen die het risico elk met \qty{5}{\%} veranderen.
\textbf{23.} Allelen binnen een blok van tientallen kilobasen worden
samen geërfd (koppelingsonevenwicht), zodat elk daarvan dezelfde
associatie vertoont; de gegenotypeerde SNP is enkel degene die op de
array stond. Om de oorzakelijke variant te vinden: sequenceer het gebied
bij gevallen en controles, versmal de verzameling statistisch, en toets
daarna elke kandidaat --- reportertoetsen voor de enhanceractiviteit van
elk allel, elk allel in cellen bewerken en de expressie van de nabije
genen meten, colokalisatie met signalen van expressiekwantitatieve
loci.
\textbf{24.} $1 - e^{-0.5} = 39\,\%$ van het genoom wordt gelezen. Het
geeft een rechtstreekse steekproef van een populatie op een bekend
tijdstip in het verleden --- allelfrequenties vóór latere migraties en
vermenging, de lengte van neanderthalersegmenten (lang, omdat weinig
generaties van recombinatie ze hadden gebroken) en daarmee een datering
van de vermenging --- die moderne genomen alleen via modellen kunnen
afleiden.
\textbf{25.} Ongeveer \num{13700} contigs in de proefassemblage;
$28\times$ per haploïd genoom, ongeveer $14$ \omterm{def:b3:genomics:ngs}{reads} per allel; zo'n $77$
nieuwe mutaties bij een kind.
\end{solution}
